\documentclass[10pt]{article}
\usepackage[T1]{fontenc}
\usepackage[utf8]{inputenc}
\usepackage[french]{babel}
\usepackage{amsmath,amssymb,mathrsfs,mathpazo}
\usepackage{geometry}
\geometry{a4paper,margin=2cm}
\usepackage{graphicx,array,tabularx,tabularray,colortbl,multirow,multicol,booktabs,enumitem,pifont,ulem,fancybox,diagbox,cancel}
\usepackage{pstricks,pstricks-add,pst-plot,pst-tree,pst-node,pst-text,pst-func,pst-math,pst-bspline,pst-3dplot,pst-eucl,pst-solides3d}
\let\pscalculate\undefined
\usepackage{pst-bezier}
\usepackage{tikz,pgfplots,tkz-tab}
\usetikzlibrary{arrows,arrows.meta,patterns,calc,positioning,trees,decorations.markings,shapes.geometric,shapes.misc,shapes.arrows,intersections,decorations.pathmorphing,decorations.pathreplacing,backgrounds,shadows,babel}
\pgfplotsset{compat=1.18}
\usepgfplotslibrary{fillbetween}
\usepackage[np]{numprint}
\usepackage{listings,makecell,pas-tableur,twoopt}
\let\barre\undefined
\usepackage[pstricks]{tablvar}
\definecolor{darkgreen}{RGB}{0,100,0}
\providecommand{\red}{\color{red}}
\providecommand{\blue}{\color{blue}}
\providecommand{\green}{\color{green}}
\providecommand{\black}{\color{black}}
\providecommand{\white}{\color{white}}
\providecommand{\gray}{\color{gray}}
\providecommand{\cyan}{\color{cyan}}
\providecommand{\magenta}{\color{magenta}}
\providecommand{\yellow}{\color{yellow}}
\providecommand{\violet}{\color{violet}}
\providecommand{\orange}{\color{orange}}
\providecommand{\euro}{\texteuro}
\providecommand{\EUR}{\texteuro}
\providecommand{\ieme}{\textsuperscript{e}}
\providecommand{\degres}{\ensuremath{^\circ}}
\providecommand{\degree}{\ensuremath{^\circ}}
\providecommand{\vv}[1]{\overrightarrow{#1}}
\providecommand{\Coord}[1]{#1}
\providecommand{\decosix}{}
\providecommand{\starredbullet}{\ensuremath{\bullet}}
\pagestyle{empty}
\setlength{\parindent}{0pt}
\usepackage{adjustbox}\DeclareUnicodeCharacter{2236}{:}\DeclareUnicodeCharacter{2190}{\ensuremath{\leftarrow}}
\providecommand{\R}{\mathbb{R}}
\providecommand{\N}{\mathbb{N}}
\providecommand{\Z}{\mathbb{Z}}
\providecommand{\Q}{\mathbb{Q}}
\providecommand{\D}{\mathbb{D}}
\providecommand{\C}{\mathbb{C}}
\providecommand{\e}{{\mathrm{e}}}
\providecommand{\vect}[1]{\overrightarrow{#1}}
\providecommand{\vectt}[1]{\overrightarrow{\mathrm{#1}}}
\providecommand{\vvt}[1]{\overrightarrow{\mathrm{#1}}}
\providecommand{\np}[2][]{{#2~#1}}
\providecommand{\vv}[1]{{\overrightarrow{#1}}}
\renewcommand{\d}{{\mathrm{d}}}
\providecommand{\degres}{{^\circ}}
\providecommand{\degre}{{^\circ}}
\providecommand{\blue}{\color{blue}}
\providecommand{\red}{\color{red}}
\providecommand{\enskip}{\quad}
\providecommand{\llbracket}{\lbrack\!\lbrack}
\providecommand{\rrbracket}{\rbrack\!\rbrack}
\providecommand{\euro}{{\text{€}}}
\providecommand{\no}{{\text{n°}}}
\providecommand{\interval}[2]{{[#1;#2]}}
\providecommand{\intervFO}[2]{{[#1;#2[}}
\providecommand{\dotfill}{\ldots}

\providecommand{\ds}{\displaystyle}
\providecommand{\textding}[1]{#1}
\providecommand{\Oij}{\(\mathrm{O};\vec{\imath},\vec{\jmath}\)}
\providecommand{\Oijk}{\(\mathrm{O};\vec{\imath},\vec{\jmath},\vec{k}\)}
\providecommand{\Ouv}{\(\mathrm{O};\vec{u},\vec{v}\)}
\providecommand{\intervFF}[2]{\ensuremath{[#1;#2]}}
\renewcommand{\euro}{\ifmmode\text{\texteuro}\else\texteuro\fi}
\DeclareUnicodeCharacter{00B0}{\ifmmode^{\circ}\else\textdegree\fi}
\DeclareUnicodeCharacter{20AC}{\ifmmode\text{\texteuro}\else\texteuro\fi}
% Énoncés : Éducation nationale. Transcription et corrigés : APMEP, Denis Vergès et contributeurs. Découpage et classement : CoopMaths, Sébastien Lozano, Éric Elter et contributeurs. Adaptation de la présentation : Maths Explorer.
% CC BY-SA : https://coopmaths.fr/www/about/
% https://coopmaths.fr/alea/static/bac/2023/tex/bac_2023_03_sujet2_metropole_2_cor.tex
\begin{document}
\noindent\textit{APMEP / CoopMaths. CC BY-SA. Présentation adaptée par Maths Explorer.}\par\medskip
\medskip

Des biologistes étudient l'évolution d'une population d'insectes dans un jardin botanique. 

Au début de l'étude la population est de \np{100000}~insectes.

Pour préserver l'équilibre du milieu naturel le nombre d'insectes ne doit pas dépasser \np{400000}.

\bigskip

\textbf{Partie A : Étude d'un premier modèle en laboratoire}

\medskip

L'observation de l'évolution de ces populations d'insectes en laboratoire, en l'absence de tout prédateur, montre que le nombre d'insectes augmente de 60\,\% chaque mois.

En tenant compte de cette observation, les biologistes modélisent l'évolution de la population d'insectes à l'aide d'une suite $\left(u_n\right)$ où, pour tout entier naturel $n$, $u_n$ modélise le nombre d'insectes, exprimé en millions, au bout de $n$ mois. 
On a donc $u_0 = 0,1$.

\medskip

\begin{enumerate}
\item 
Augmenter de 60\,\%, c'est multiplier par $1+\dfrac{60}{100}=1,6$.
La suite $(u_n)$ est donc géométrique de raison $q=1,6$.

La forme explicite d'une suite géométrique de raison $q$ et de premier terme $u_0$ est: $u_n=u_0\times q^n$ donc $u_n = 0,1 \times 1,6^n$ pour tout $n$ de $\N$.

\item 
$1,6>1$ donc $\ds\lim_{n\to +\infty} 1,6^n = +\infty$

On en déduit que $\ds\lim_{n\to +\infty} 0,1\times 1,6^n = +\infty$ et donc que $\ds\lim_{n\to +\infty} u_n = +\infty$.

\item On résout l'inéquation $u_n > 0,4$.

$\begin{gathered}u_n>0,4
\iff
0,1\times 1,6^n >0,4
\iff
1,6^n >4
\iff
\ln\left (1,6^n\right ) > \ln(4)
\iff
n\times \ln(1,6)>\ln(4)\\
\phantom{u_n>0,4}
\iff
n>\dfrac{\ln(4)}{\ln(1,6)}\end{gathered}$

Or $\dfrac{\ln(4)}{\ln(1,6)}\approx 2,95$, donc le plus petit entier naturel $n$ à partir duquel $u_n>0,4$ est 3.

\item 
$u_3>0,4$ signifie que le nombre d'insectes dépasse \np{400000} dès le 3\ieme{} mois; selon ce modèle le milieu naturel n'est donc pas préservé.

\end{enumerate}

\bigskip

\textbf{Partie B : Étude d'un second modèle}

\medskip

En tenant compte des contraintes du milieu naturel dans lequel évoluent les insectes, les biologistes choisissent une nouvelle modélisation.
Ils modélisent le nombre d'insectes à l'aide de la suite $\left(v_n\right)$, définie par :
$v_0 = 0,1$  et, pour tout entier naturel $n,\: v_{n+1} = 1,6v_n - 1,6v_n^2$,
où, pour tout entier naturel $n,\: v_n$ est le nombre d'insectes, exprimé en millions, au bout de $n$ mois.

\medskip

\begin{enumerate}
\item $v_1=1,6v_0 - 1,6 v_0^2 = 1,6\times 0,1 - 1,6\times 0,1^2 = 0,144$.

Le nombre d'insectes au bout d'un mois est donc égal à \np{144000}.

\item On considère la fonction $f$ définie sur l'intervalle $\left[0~;~\dfrac12\right]$ par
$f(x) = 1,6x - 1,6x^2$.

	\begin{enumerate}
		\item On résout l'équation $f(x) = x$.

$\begin{gathered}f(x)=x
\iff
1,6x-1,6x^2=x
\iff
0,6x - 1,6x^2=0
\iff
x\left (0,6-1,6x\right )=0\\
\phantom{f(x)=x}
\iff
x=0 \text{ ou }	0,6-1,6x=0
\iff
x=0 \text{ ou } x=\dfrac{0,6}{1,6}
\iff
x=0 \text{ ou } x=\dfrac{3}{8}\end{gathered}$

Les deux solutions appartiennent à l'intervalle $\left[0~;~\dfrac12\right]$, donc l'équation $f(x)=x$ admet deux solutions dans cet intervalle: 0 et $\dfrac{3}{8}$.
		
		\item 
$f'(x)=1,6\times 1 - 1,6\times 2x = 1,6\left (1-2x\right )$

$x\in\left[0~;~\dfrac12\right]$ donc $x\leqslant \dfrac{1}{2}$ et donc $1-2x\geqslant 0$.

Sur $\left[0~;~\dfrac12\right]$, $f'(x)\geqslant 0$ donc $f$ est croissante.
			\end{enumerate}
\item 
	\begin{enumerate}
		\item On va montrer par récurrence que la propriété $n$, $0 \leqslant v_n  \leqslant  v_{n+1} \leqslant  \dfrac12$ est vraie pour tout entier naturel $n$.
		
\begin{list}{\textbullet}{}
\item \textbf{Initialisation}

$v_0 = 0,1$ et $v_1=\dfrac{3}{8}$; on a donc $0 \leqslant v_0 \leqslant v_1 \leqslant \dfrac{1}{2}$.

La propriété est donc vraie au rang 0.
\item \textbf{Hérédité}

On suppose la propriété vraie au rang $n$, c'est-à-dire $0 \leqslant v_n  \leqslant v_{n+1} \leqslant  \dfrac12$; c'est l'hypothèse de récurrence.

On a: $0 \leqslant v_n  \leqslant  v_{n+1} \leqslant  \dfrac12$ et on sait que la fonction $f$ est croissante sur $\left[0~;~\dfrac12\right]$; on en déduit que:
$f(0) \leqslant f(v_n)  \leqslant  f(v_{n+1}) \leqslant  f\left (\dfrac12\right )$.

$f(0)=0$; $f\left(v_n\right)=v_{n+1}$; $f\left(v_{n+1}\right)=v_{n+2}$ et $f\left (\dfrac{1}{2}\right )=1,6\times \dfrac{1}{2}-1,6\times \left (\dfrac{1}{2}\right )2=0,4$

Donc $f(0) \leqslant f(v_n)  \leqslant  f(v_{n+1}) \leqslant  f\left(\dfrac12\right )$ équivaut à $0 \leqslant v_{n+1}  \leqslant v_{n+2} \leqslant 0,4$
ce qui entraine $0 \leqslant v_{n+1} \leqslant v_{n+2} \leqslant \dfrac{1}{2}$.

La propriété est donc vraie au rang $n+1$.
\item \textbf{Conclusion}

La propriété est vraie au rang 0 et elle est héréditaire pour tout $n \in \N$; d'après le principe de récurrence, elle est vraie pour tout $n \in \N$.
\end{list}

On a donc démontré que, pour tout $n$, on a: $0 \leqslant v_n  \leqslant v_{n+1} \leqslant  \dfrac12$.
		
		\item 
\begin{list}{\textbullet}{On sait que:}
\item $v_n \leqslant v_{n+1}$ pour tout $n$, donc la suite $\left(v_n\right)$ est croissante;
\item $v_n\leqslant \dfrac{1}{2}$ pour tout $n$, donc la suite $\left(v_n\right)$ est majorée par $\dfrac{1}{2}$.
\end{list}

La suite $\left(v_n\right)$ est croissante et majorée donc, d'après le théorème de la convergence monotone, on peut dire que la suite $\left(v_n\right)$ est convergente.

On note $\ell$ la valeur de sa limite. On admet que $\ell$ est solution de l'équation $f(x) = x$.

		\item 
La suite $(v_n)$ est croissante et admet pour limite $\ell$; donc pour tout $n$, on aura $v_n \leqslant \ell$. En particulier $v_1\leqslant \ell$ donc $\ell\geqslant 0,1$.

La limite $\ell$ est solution de l'équation $f(x) = x$ donc c'est soit $0$, soit $\dfrac{3}{8}$. Mais $\ell\geqslant 0,1$ donc $\ell$ ne peut être égale à 0.
Donc $\ell=\dfrac{3}{8}=0,375$.		
		
Pour tout $n$, on aura $v_n\leqslant \ell$, donc $v_n \leqslant 0,375$; il y aura donc toujours moins de \np{375000} insectes. Donc, selon ce modèle,  l'équilibre du milieu naturel sera préservé.
	\end{enumerate}	
\end{enumerate}

\begin{enumerate}[resume]
\item  On donne ci-contre la fonction \texttt{seuil}, écrite en langage Python.

\begin{minipage}{0.6\linewidth}
	\begin{enumerate}
		\item 
La fonction \texttt{seuil(a)} donne la première (et plus petite) valeur de \texttt{n} telle que \texttt{v>=a}, c'est-à-dire telle que $v_n\geqslant a$.

On a vu que $v_n\leqslant 0,375$ pour tout $n$; il n'y a donc pas de valeur de $n$ pour laquelle $v_n\geqslant 0,4$.

Le programme ne s'arrête jamais.
	\end{enumerate}
\end{minipage}\hfill
\begin{minipage}{0.36\linewidth}
\begin{tabularx}{\linewidth}{|X|}\hline
def seuil(a) :\\
\qquad v=0.1\\
\qquad n=0\\
\qquad while v<a :\\
\qquad \qquad v=1.6*v-1.6*v*v\\
\qquad \qquad n=n+1\\
\qquad return n\\ \hline
\end{tabularx}
\end{minipage}

\begin{enumerate}[resume]
		\item 
À la calculatrice, on trouve $v_5\approx 0,338 <0,35$ et $v_6\approx 0,358 \geqslant 0,35$; donc la valeur renvoyée par \texttt{seuil(0.35)} est 6.
		

Cela signifie qu'à partir du 6\ieme{} mois, il y aura plus de \np{350000} insectes.
\end{enumerate}	

\end{enumerate}	

\bigskip

\end{document}
